\begin{frame}{직관: ``이 상태의 가격표''}
  \begin{itemize}
    \item $V^\pi$ 는 상태 공간에 \textbf{매 상태마다 하나의 숫자}를
      매기는 함수 --- 상태별 ``\kb{가격표}''다.
    \item $Q^\pi$ 는 ``상태-행동'' \textbf{조합}마다 하나의 숫자를
      매기는, 좀 더 \textbf{잘게 쪼갠} 가격표.
    \item ``이 상태에서 앞으로 평균적으로 얼마를 벌 수 있나''라는
      질문에 $V^\pi(s)$ 가 바로 답을 주는 숫자.
  \end{itemize}
  \vskip0.5em
  \begin{block}{중요: ``실제 리턴''이 아니라 ``기댓값''}
    이 숫자는 ``\textbf{실제로 벌어지는 리턴}''이 아니라 그
    \textbf{기댓값}이다. 같은 상태에서 시작해도 전이가 확률적이면
    매 에피소드마다 리턴이 다르다 --- $V^\pi(s)$ 는 ``이 상태에서
    출발하면 \textbf{평균적으로} 이 정도''를 말해주는 \kb{요약
    통계}다.
  \end{block}
\end{frame}
